\documentclass[11pt,a4paper]{article}

\usepackage{fontspec}
\usepackage{polyglossia}
\setmainlanguage{vietnamese}
\setotherlanguage{english}
\setmainfont{Times New Roman}
\setsansfont{Helvetica Neue}
\setmonofont{Menlo}[Scale=0.82]

\usepackage[margin=2.4cm,top=2.6cm,bottom=2.4cm]{geometry}
\usepackage{booktabs}
\usepackage{amsmath,amssymb}
\usepackage{xcolor}
\usepackage{listings}
\usepackage{titlesec}
\usepackage{enumitem}
\usepackage{fancyhdr}
\usepackage[hidelinks]{hyperref}

\definecolor{hcmusblue}{HTML}{1B3A5C}
\definecolor{contextblue}{HTML}{2E86C1}
\definecolor{codebg}{HTML}{F5F7FA}
\definecolor{coderule}{HTML}{2E5C8A}
\definecolor{codetext}{HTML}{A03030}

\titleformat{\section}{\sffamily\large\bfseries\color{hcmusblue}}
            {\thesection.}{0.6em}{}[\vspace{2pt}\hrule]
\titleformat{\subsection}{\sffamily\normalsize\bfseries\color{contextblue}}
            {\thesubsection}{0.6em}{}
\titleformat{\subsubsection}{\sffamily\small\bfseries}
            {\thesubsubsection}{0.6em}{}
\titlespacing*{\section}{0pt}{14pt}{6pt}
\titlespacing*{\subsection}{0pt}{10pt}{4pt}

\setlist[itemize]{leftmargin=1.3em, topsep=3pt, itemsep=2pt}
\setlist[enumerate]{leftmargin=1.5em, topsep=3pt, itemsep=2pt}

\newcommand{\code}[1]{{\ttfamily\small\color{codetext}#1}}

\lstset{
  basicstyle=\ttfamily\scriptsize,
  backgroundcolor=\color{codebg},
  frame=leftline, framerule=1.5pt, rulecolor=\color{coderule},
  framesep=6pt, xleftmargin=8pt,
  breaklines=true, breakatwhitespace=false,
  showstringspaces=false, columns=fullflexible,
}

\pagestyle{fancy}
\fancyhf{}
\fancyhead[L]{\footnotesize\color{gray}APKGC --- Báo cáo chuyên đề}
\fancyhead[R]{\footnotesize\color{gray}\thepage}
\renewcommand{\headrulewidth}{0.4pt}

\begin{document}

\begin{center}
  {\sffamily\LARGE\bfseries\color{hcmusblue} BÁO CÁO CHUYÊN ĐỀ}\\[8pt]
  {\sffamily\large\bfseries APKGC: Noise-enhanced Multi-Modal Knowledge Graph\\
   Completion with Attention Penalty}\\[10pt]
  {\small Seminar GraphML --- Học phần HP3}\\[2pt]
  {\small ĐHQG-HCM --- Trường Đại học Khoa học Tự nhiên}
\end{center}

\vspace{4pt}
\begin{center}\small
\begin{tabular}{@{}ll@{}}
  \toprule
  \textbf{Hội nghị} & The 39th AAAI Conference on Artificial Intelligence
                      (AAAI-25), Oral, tr. 15005--15013\\
  \textbf{Tác giả}  & Yue Jian, Xiangyu Luo, Zhifei Li, Miao Zhang, Yan Zhang,
                      Kui Xiao, Xiaoju Hou\\
                    & Hubei University, Trung Quốc\\
  \textbf{Mã nguồn} & \texttt{github.com/HubuKG/APKGC}\\
  \bottomrule
\end{tabular}
\end{center}

\vspace{6pt}
\section{Dàn ý trình bày --- ImparTail}
\textbf{Bài báo:} Towards Unbiased and Robust Spatio-Temporal Scene Graph Generation

and Anticipation --- Peddi, Saurabh, Shrivastava, Singla, Gogate. CVPR 2025 (Highlight).

\textbf{Bộ slide:} \code{slides/ImparTail\_Slides.pdf}, 26 trang. \textbf{Thời lượng:} 60 phút.

\section{Phân bổ thời gian}
\begin{center}\small
\begin{tabular}{@{}llrr@{}}
\toprule
\textbf{Phần} & \textbf{Slide} & \textbf{Phút} & \textbf{Cộng dồn}\\
\midrule
1. Mở đầu & 1--2 & 2 & 2\\
2. Kiến thức nền & 3--5 & 9 & 11\\
3. Vấn đề đuôi dài & 6--7 & 7 & 18\\
4. Phương pháp & 8--12 & 17 & 35\\
5. Đóng góp thứ ba (độ bền) & 13 & 6 & 41\\
6. Thực nghiệm & 14--16 & 8 & 49\\
7. Đọc phản biện & 17--19 & 8 & 57\\
8. Kết \& hỏi đáp & 20 & 3 & 60\\
\bottomrule
\end{tabular}
\end{center}
Nếu cháy giờ: cắt slide 7 (các hướng đã có) và rút gọn slide 16 xuống 2 phút. \textbf{Không cắt} slide 12 (bộ ba đi xuyên mô hình) và 17 (phát hiện về curriculum) --- đó là hai chỗ người nghe nhớ nhất.

\section{Phần 1 --- Mở đầu (2 phút)}
Mở bằng một câu hỏi, không phải bằng tên bài báo:

> "Nhìn một đoạn video người ta pha cà phê. Máy tính cần trả lời được: ai đang > làm gì với cái gì, và \textbf{sắp} làm gì tiếp. Đó là bài toán hôm nay."

Nói rõ đây là \textbf{CVPR 2025 Highlight} --- top khoảng 3\% số bài nộp. Rồi vào luôn.

\section{Phần 2 --- Kiến thức nền (9 phút)}
\subsection{Slide 3 --- Scene graph: ảnh tĩnh sang video (3 phút)}
Bắt đầu từ cái quen thuộc: scene graph tĩnh là đồ thị, đỉnh là vật thể, cạnh là quan hệ. Một câu về Visual Genome.

Rồi chỉ vào hình ba khung: \textbf{cùng một cặp (người, cốc)}, quan hệ đổi từ \emph{nhìn} $\rightarrow$ \emph{chạm} $\rightarrow$ \emph{uống}.

> "Đây là thứ scene graph tĩnh không diễn tả được. Ba tấm ảnh rời cho ba đồ thị > rời. Cái ta cần là \textbf{một chuỗi đồ thị có liên hệ với nhau theo thời gian}."

\subsection{Slide 4 --- Ký hiệu (3 phút)}
Đi chậm, vì cả bài dùng lại. Viết lên bảng nếu có:

\begin{itemize}\itemsep2pt
\item Đoạn video \$V\_1\textasciicircum{}T\$ gồm \$T\$ khung
\item Trong khung \$t\$: vật thể \$o\_k\textasciicircum{}t\$, mỗi cái có hộp bao và nhãn lớp
\item Vị ngữ \$p\_\{ijk\}\textasciicircum{}t\$ nối cặp \$(o\_i\textasciicircum{}t, o\_j\textasciicircum{}t)\$
\item \textbf{Bộ ba} \$r\_\{ijk\}\textasciicircum{}t = (o\_i\textasciicircum{}t, p\_\{ijk\}\textasciicircum{}t, o\_j\textasciicircum{}t)\$ --- đây là đơn vị cơ bản
\end{itemize}
Rồi hai tác vụ, nhấn vào chỗ khác nhau:

> "VidSGG nhìn khung nào thì mô tả khung đó. SGA phải \textbf{đoán khung chưa nhìn > thấy}. Cùng bộ máy, nhưng SGA khó hơn hẳn."

\subsection{Slide 5 --- Ví dụ đầu vào (3 phút)}
Đây là slide đầu tiên người nghe thực sự \emph{thấy} bài toán. Đi từ trái sang phải:

\begin{enumerate}\itemsep2pt
\item Hai khung đã quan sát $\rightarrow$ VidSGG sinh (người, \textbf{looking\_at}, sổ tay)
\item Mũi tên SGA
\item Hai khung chưa quan sát $\rightarrow$ dự đoán (người, \textbf{touching}, cốc) rồi
\end{enumerate}
(người, \textbf{drinking\_from}, cốc)

> "Chú ý chuỗi \emph{chạm} rồi \emph{uống}. Mô hình phải học rằng chạm vào cốc thường dẫn > tới uống. Đó là tri thức về \textbf{diễn tiến}, không phải về một khung hình."

\section{Phần 3 --- Vấn đề đuôi dài (7 phút)}
\subsection{Slide 6 --- Phân phối đuôi dài (4 phút)}
Chỉ vào biểu đồ cột. Bốn cột xanh cao vống, phần đuôi đỏ dẹt lét.

Câu chốt của cả phần này:

> "Vấn đề không phải là đuôi dài. Vấn đề là \textbf{lớp ở đuôi mới là lớp đáng giá}. > Biết 'người ở phía trước cái cốc' gần như vô dụng. Biết 'người đang uống từ > cái cốc' mới là hiểu cảnh. Mà mô hình lại học ngược lại."

Giải thích vì sao \textbf{mean Recall} chứ không Recall:

> "Recall@K cộng gộp mọi dự đoán, nên lớp head chiếm đa số và che mất điểm yếu > ở đuôi. mean Recall tính recall \textbf{từng lớp rồi lấy trung bình} --- mỗi lớp một > phiếu. Đổi chỉ số là đổi cả bài toán."

\subsection{Slide 7 --- Các hướng đã có (3 phút) --- \emph{cắt được nếu thiếu giờ}}
Lướt nhanh bảng bốn hướng. Dừng ở dòng cuối:

> "Ba hướng trên đều \textbf{thêm tham số hoặc thêm mô-đun}. ImparTail không thêm gì > cả --- nó chỉ đổi hàm mất mát. Đó là lý do nó gắn được lên sáu mô hình nền khác > nhau."

\section{Phần 4 --- Phương pháp (17 phút, trọng tâm)}
\subsection{Slide 8 --- Toàn cảnh pipeline (3 phút)}
Chỉ vào dấu ngoặc lớn: \textbf{ba khối đầu lấy nguyên từ mô hình nền}. Chỉ khối cuối là của ImparTail.

> "Nói cách khác: forward pass \textbf{không đổi một dòng nào}. Chỉ backward đổi."

\subsection{Slide 9 --- Masked Loss (5 phút, chậm nhất phần này)}
Viết công thức lên bảng:

\$\$\textbackslash\{\}widetilde\{\textbackslash\{\}mathcal\{L\}\}\_\{p\} = (1 - m) \textbackslash\{\}ast \textbackslash\{\}mathcal\{L\}\_\{p\}\$\$

Rồi giải thích bằng lời trước khi giải thích bằng ký hiệu:

> "\$m\$ là một cái công tắc nhị phân cho \textbf{từng bộ ba một}. Bật lên thì mất mát > của bộ ba đó bằng không --- nó biến mất khỏi gradient, như thể không tồn tại > trong batch này."

Nhấn chỗ khác biệt với cost-sensitive learning:

> "Đánh trọng số lại là \emph{giảm tiếng nói} của lớp head. Che là \textbf{tắt hẳn}. Nghe > thô, nhưng chính vì thô mà không phải dò tham số trọng số."

Chỉ vào hình bên phải: hai thanh xám gạch đứt là head bị che, hai thanh đỏ là tail còn nguyên.

\subsection{Slide 10 --- Hàm mục tiêu (3 phút)}
Đừng đọc từng số hạng. Chỉ nói hai ý:

\begin{enumerate}\itemsep2pt
\item VidSGG có hai số hạng: vật thể và vị ngữ. \textbf{Chỉ vị ngữ bị che.}
\item SGA có năm số hạng. \textbf{Hai chỗ bị che}: vị ngữ ở khung đã quan sát và vị ngữ
\end{enumerate}
ở khung dự đoán. Hộp bao và tái tạo giữ nguyên.

> "Vì sao chỉ che nhánh vị ngữ? Vì chỉ nhánh đó mới bị đuôi dài. Phân bố lớp vật > thể cân bằng hơn nhiều."

\subsection{Slide 11 --- Curriculum (3 phút)}
Ba thành phần: xếp hạng theo lượng dữ liệu, đo thành thục bằng recall, trộn dần theo epoch. Tỉ lệ che \$\textbackslash\{\}mathcal\{R\}\_m = e \textbackslash\{\}times \textbackslash\{\}mathcal\{R\}\_s\$ tăng dần.

Chỉ vào ba biểu đồ bên phải: phân phối nhãn còn lại dịch từ đuôi dài sang đều.

\textbf{Gieo mầm cho slide 17} --- đừng nói hết ngay:

> "Ghi nhớ slide này. Lát nữa ta sẽ quay lại xem nó có thực sự được dùng không."

\subsection{Slide 12 --- Một bộ ba đi xuyên mô hình (3 phút)}
Đi từng bước với một bộ ba cụ thể:

\begin{enumerate}\itemsep2pt
\item Khung \$I\textasciicircum{}\{12\}\$ có người và cốc
\item ORPU cho đặc trưng, hộp bao, phân phối lớp
\item STPU cho biểu diễn quan hệ \$z\_\{ij\}\textasciicircum{}\{12\}\$
\item Decoder cho phân phối vị ngữ
\item Nhãn thật là \textbf{looking\_at} $\rightarrow$ thuộc head $\rightarrow$ \$m = 1\$ $\rightarrow$ \textbf{mất mát bằng 0}
\end{enumerate}
> "Bộ ba này vừa đi hết mô hình rồi \textbf{không đóng góp gì cho việc học}. Nếu nhãn > là \emph{drinking\_from} thì ngược lại --- gradient chảy qua bình thường."

\section{Phần 5 --- Đóng góp thứ ba (6 phút)}
\subsection{Slide 13 --- Robust VidSGG và Robust SGA (6 phút)}
Giải thích thiết kế: \textbf{huấn luyện trên dữ liệu sạch, đánh giá trên dữ liệu nhiễu}. Không phải huấn luyện tăng cường.

> "Đây mô phỏng đúng tình huống triển khai thật: mô hình được huấn luyện trong > phòng, rồi đem ra chỗ có sương mù, chói nắng, camera nhiễu."

16 kiểu nhiễu, ba mức. Nhóm lại cho dễ nhớ: nhiễu cảm biến, thời tiết, ánh sáng.

Nhấn: \textbf{đây là hai tác vụ mới bài báo đặt ra}, không phải chỉ là thí nghiệm phụ.

\section{Phần 6 --- Thực nghiệm (8 phút)}
\subsection{Slide 14 --- Thiết lập (2 phút)}
Nói cách đọc bảng trước khi chiếu bảng:

> "Mỗi dòng là một mô hình nền, dòng ngay dưới là \textbf{chính mô hình đó} sau khi > gắn ImparTail. So theo cặp, đừng so chéo giữa các mô hình."

Ba chế độ VidSGG khác nhau ở lượng thông tin cho sẵn: SGDET chỉ có khung hình, SGCLS thêm hộp bao, PREDCLS thêm cả nhãn vật thể.

\subsection{Slide 15 --- Kết quả VidSGG (3 phút)}
Đừng đọc số. Chỉ vào các ô xanh:

> "Cải thiện ở \textbf{mọi ô}, không có ngoại lệ. Mức tăng mR@10 khoảng 12\% cho > SGDET, 22\% cho SGCLS, trên 20\% cho PREDCLS."

Nếu có người hỏi vì sao PREDCLS tăng nhiều nhất: vì cho sẵn nhãn vật thể nên phần khó còn lại đúng là phân loại vị ngữ --- chính chỗ ImparTail tác động.

\subsection{Slide 16 --- SGA và độ bền (3 phút)}
Bảng SGA: khoảng 12\% cho STTran++/DSGDetr++, 20\% cho SceneSayer.

Biểu đồ độ bền là chỗ đáng dừng:

> "Cả hai đều tụt khi gặp nhiễu --- đó là bình thường. Điều đáng nói là mô hình > ImparTail \textbf{tụt ít hơn}: 8,0 điểm so với 14,3. Và điểm sau khi nhiễu của nó > vẫn \textbf{cao hơn điểm sạch} của mô hình nền."

Giải thích cơ chế theo bài báo: cân bằng phân phối làm mô hình không quá khớp vào mẫu hình phổ biến, nên nhiễu ít phá hơn --- giống tác dụng của dropout.

\section{Phần 7 --- Đọc phản biện (8 phút)}
\subsection{Slide 17 --- Curriculum không được dùng để lấy số (4 phút)}
Đây là điểm nhấn của cả bài nói. Dựng lên từ từ.

Nhắc lại slide 11, rồi chiếu nguyên văn Remark trang 7. Đọc to đoạn \emph{"we opted for a fixed Masking Ratio of 0.9"}.

Rồi đưa ba tầng bằng chứng:

\begin{center}\small
\begin{tabular}{@{}ll@{}}
\toprule
\textbf{Nguồn} & \textbf{Bằng chứng}\\
\midrule
Bài báo & Remark trang 7 nói rõ dùng tỉ lệ cố định\\
Mã nguồn & \code{train\_sgg\_methods.py:402} $\rightarrow$ \code{is\_curriculum = False} là mặc định\\
Script chạy & \code{METHODS=("sttran" "dsgdetr")} --- không script nào gọi curriculum\\
\bottomrule
\end{tabular}
\end{center}
\textbf{Phát biểu cẩn thận} --- đây là chỗ dễ nói quá:

> "Tôi chỉ đang mô tả \textbf{bài báo và mã đã công bố}. Không kết luận tác giả có > chạy curriculum hay không trong quá trình nghiên cứu --- rất có thể họ có chạy > mà chỉ không đưa vào bản phát hành."

Câu hỏi để lại:

> "Nếu tỉ lệ cố định cho kết quả tương đương curriculum, thì đóng góp thật nằm ở > curriculum hay chỉ ở việc che mất mát? Bài báo chưa tách được hai thứ đó."

\subsection{Slide 18 --- Quan sát khác khi đọc mã (2 phút)}
Lướt nhanh: đường dẫn cứng của bốn máy, không cố định seed, rào cản chạy lại.

Nói rõ vì sao nhóm không chạy lại được --- và \textbf{không đổ cho mã hỏng}:

> "Không phải mã sai. Là chi phí dữ liệu: hàng chục GB video Charades, bước tách > khung hình, và \code{torch 1.10.1+cu111} --- bản năm 2021 mà GPU hiện đại không còn > hỗ trợ."

Ghi nhận điểm cộng: giấy phép Apache 2.0 rõ ràng, có phát hành checkpoint --- khá hơn nhiều repo cùng lĩnh vực.

\subsection{Slide 19 --- Mạnh yếu (2 phút)}
Đọc lướt hai cột. Dừng ở dòng "tỉ lệ che 0,9":

> "Che 90\% nhãn mà vẫn học được --- con số này tự nó đã đáng ngạc nhiên. Bài báo > không khảo sát vì sao lại là 0,9."

\section{Phần 8 --- Kết (3 phút)}
\subsection{Slide 20 --- Ba điều mang về}
Nói chậm, mỗi ý một câu:

\begin{enumerate}\itemsep2pt
\item \textbf{Đổi mục tiêu, đừng đổi kiến trúc.} Một phép nhân với mặt nạ nhị phân, không
\end{enumerate}
thêm tham số nào, cải thiện 12--22\%.

\begin{enumerate}\itemsep2pt
\item \textbf{Chỉ số quyết định điều mình nhìn thấy.} Recall giấu điểm yếu ở đuôi, mean
\end{enumerate}
Recall phơi ra.

\begin{enumerate}\itemsep2pt
\item \textbf{Đọc bài báo cùng với mã.} Phát hiện về curriculum chỉ hiện ra khi đọc cả
\end{enumerate}
hai; đọc riêng cái nào cũng không thấy.

\section{Câu hỏi có thể bị hỏi}
\textbf{1. Che 90\% nhãn thì còn gì để học?}

Che theo \textbf{bộ ba}, không theo mẫu. Một khung hình có nhiều cặp vật thể; che hết các cặp head vẫn còn các cặp tail. Và mất mát vật thể không bị che.

\textbf{2. Khác gì focal loss?}

Focal loss giảm trọng số mẫu dễ theo \textbf{độ tin cậy dự đoán}, liên tục và tự động. ImparTail che theo \textbf{lớp}, nhị phân, và theo lịch trình định trước. Bài báo không so trực tiếp với focal loss --- đó là một thiếu sót.

\textbf{3. Vì sao không huấn luyện tăng cường với nhiễu cho robust?}

Vì tác vụ họ đặt ra cố tình mô phỏng triển khai thật: không biết trước sẽ gặp nhiễu gì. Huấn luyện tăng cường sẽ trả lời một câu hỏi khác.

\textbf{4. Cải thiện có nằm trong nhiễu ngẫu nhiên không?}

\textbf{Không kiểm chứng được} --- bài báo không báo cáo độ lệch chuẩn, không nói chạy

mấy seed. Với mức tăng 12--22\% thì khả năng là thật, nhưng đúng ra phải có nhiều seed.

\textbf{5. Nhóm có chạy lại không?}

Không, và lý do là chi phí dữ liệu chứ không phải mã hỏng --- xem slide 18. Nhóm có chạy lại đầy đủ bài APKGC trong phần báo cáo.

\textbf{6. ImparTail có gắn được lên mô hình mới hơn không?}

Về nguyên tắc có, miễn mô hình đó dùng pipeline biểu diễn hướng-đối-tượng và có một số hạng mất mát phân loại vị ngữ tách rời. Bài báo đã chứng minh trên sáu mô hình nền.

\textbf{7. Vì sao chọn Action Genome?}

Đây là bộ chuẩn de-facto cho VidSGG và SGA, chú giải trên video Charades. Nhưng \textbf{chỉ một bộ dữ liệu} cũng là điểm yếu --- không biết kết quả có chuyển sang bộ khác không.

\end{document}
